\documentclass[12pt]{ctexart}
\usepackage[utf8]{inputenc}
\usepackage{geometry}
\geometry{a4paper, margin=1in}
\usepackage{hyperref}
\usepackage{listings}
\usepackage{xcolor}
\usepackage{fontspec}
\usepackage{fancyhdr}


\definecolor{codegreen}{rgb}{0,0.6,0}
\definecolor{codegray}{rgb}{0.5,0.5,0.5}
\definecolor{codepurple}{rgb}{0.58,0,0.82}
\definecolor{backcolour}{rgb}{0.95,0.95,0.92}

\lstdefinestyle{mystyle}{
    backgroundcolor=\color{backcolour},   
    commentstyle=\color{codegreen},
    keywordstyle=\color{magenta},
    numberstyle=\tiny\color{codegray},
    stringstyle=\color{codepurple},
    basicstyle=\ttfamily\footnotesize,
    breakatwhitespace=false,         
    breaklines=true,                 
    captionpos=b,                    
    keepspaces=true,                 
    numbers=left,                    
    numbersep=5pt,                  
    showspaces=false,                
    showstringspaces=false,
    showtabs=false,                  
    tabsize=2
}

\lstset{style=mystyle}

\title{半导体物理}
\author{胥宇龙}
\date{\today}
\begin{document}
    
\maketitle

\tableofcontents
\newpage

\section{晶体结构}
\subsection{原胞}
\subsection{晶胞}
\subsubsection{简单立方SC}
立方体的各个顶点有一个原子
\subsubsection{体心立方BCC}
立方体的各个顶点及体心有一个原子
\subsubsection{面心立方FCC}
立方体的各个顶点及面心有一个原子，此同时为最密堆积
\subsubsection{金刚石结构}
双面心立方晶胞，沿体对角线平移四分之一长度，两种原子组成的两套晶胞按此方法排列则为闪锌矿结构

金刚石结构的解理面为\{111\}面心立方，闪锌矿结构的解理面为\{110\}面，因不同原子间有附加的库伦作用，使\{111\}面间联系加强
\subsection{晶格结构描述}
\subsubsection{晶向指数}
\subsubsection{晶面指数}

\subsection{原子价键}
\subsection{固体中的缺陷与杂质}
对于半导体，其晶体中的缺陷将极大影响半导体器件的物理与化学性质
\subsubsection{点缺陷}
分为肖特基缺陷语与弗伦克尔缺陷，其中肖特基缺陷中仅有空位产生，弗伦克尔缺陷则同时产生空位与填隙，对于空位，其产生受主(acceptor)作用，填隙则为施主(doner)作用

对于化合物，其同时可存在反结构缺陷，即化合物中的一种原子被另一种原子替代

\subsubsection{线缺陷}
位错为晶体中常见的线缺陷，位错原子相比周围原子缺少成键，其可同时作为施主与受主存在
\subsubsection{面缺陷}
层错为主要的面缺陷情况，即其中晶体中面不连续情况的出现
\subsubsection{杂质}
晶体中与本体原子不同的原子为杂质原子，其可为人为加入或无意加入。

其分为间隙式与替位式，间隙式杂质的原子的大小较小，替位式则较大，且处于此状态时，杂质容易激活

向晶体中引入杂质的方法称为掺杂，常见的方式有高温扩散掺杂（常见温度为700-1000摄氏度）
\subsection{半导体材料的生长}
\subsubsection{熔体生长（切克劳斯基法）}
将籽晶放于熔融状态的物质中提拉生长
\subsubsection{外延生长法}
在单晶衬底上生长一层薄单晶，可分为同质外延与异质外延
化学气相沉积：
气相反应物于晶体表面发生反应，生成一层薄单晶
液相沉积：

分子束外延
\section{固体量子论初步}
\subsection{允带与禁带}
\subsubsection{电子的共有化运动}
原子组成晶体后，由于壳层的交叠，电子不再局限于一个原子上，可在相似壳层上进行转移
\subsubsection{允带与禁带的形成}
当电子进行公有化后，原子的能级将分裂为能带
\end{document}
